\section{Model and method}
\label{model}
\subsection{Pseudo-spin Model}
Following the earlier model~\cite{ishizuka2011quantum,vijigiri2020dipole}, we consider a pseudo-spin model where the displacement of protons between the $O-O$ double-well potential is mapped on to the $z$-component ($\sigma^z_i$) of a spin-1/2 operator, $\sigma_i$, i.e, $\sigma^z_i=\pm 1$. Here $+1$ $(-1)$ refers to the proton being far from (near to) the center of the $O-O$ bond~\cite{blinc1960isotopic,horiuchi2008organic}. At zeroth order we consider a  four-spin interaction (with strength $J_0$)  as discussed in earlier studies~\cite{wesselinowa1995green} and has some intricate connection to the vertex-based models~\cite{stilck1981vertex}. The physical origin of the four-spin interaction term is primarily due to the fact that it yields the minimum energy configuration mediated by the $\pi$-electrons in $\text{C}_4\text{O}_4$ by the rearrangement of $\pi$-bonds. The model Hamiltonian also includes a next-to-next nearest neighbor (NNN) Ising interaction as well as a dipole-dipole interaction signified by the coupling strength $J_1$ and $J_2$ respectively. In addition to these terms~\cite{chern2014gauge}, in the present study we introduce an additional coupling between the two-dimensional layers represented by $J_3$ . The full Hamiltonian is thus given by,
\begin{eqnarray}
\label{Hamiltonian}
\mathcal{H} &=& -J_0\sum_{\square} \hat{A}_{\square} + J_1\sum_{\square}\hat{B}_{\square}
 - J_2\sum_{\langle AB \rangle}\hat{\vec{P}}_{A}\cdot\hat{\vec{P}}_{B} \nonumber  \\
 && + J_3\sum_{L, i}\sigma_{L,i}^{z}\sigma_{L+1,i}^{z} - K\sum_{i}\sigma_{i}^{x} . %\\
 %\mathcal{H} &=& \mathcal{H}_{\text{2D}} + J_3\sum_{L, i}\sigma_{L,i}^{z}\sigma_{L+1,i}^{z}
\end{eqnarray}
In the above $\hat{A}_{\square} = \sigma_{1}^{z}\sigma_{2}^{z}\sigma_{3}^{z}\sigma_{4}^{z}$, the indices $1$,$2$,$3$,$4$ correspond to the indices of four spins on the links (green lines in Fig.~\ref{lattice}) of a given unit-plaquette ($\square$) taken counterclockwise starting from the lower horizontal link($-$) (see Fig.~\ref{lattice}). The indices 1, 3, and 2, 4 are designated to be diagonally opposite to each other.  Similarly, $\hat{B}_{\square}=(\sigma_{1}^{z}\sigma_{3}^{z} + \sigma_{2}^{z}\sigma_{4}^{z})$. $\hat{P}_{A,B}^x=\pm\frac{1}{4}(\sigma^z_{1}+\sigma^z_{4}-\sigma^z_{2}-\sigma^z_{3})$, $\hat{P}_{A,B}^y=\pm\frac{1}{4}(\sigma^z_{1}+\sigma^z_{2}-\sigma^z_{3}-\sigma^z_{4})$ are the dipole-moment vectors for A, B sub-plaquettes as illustrated in Fig.~\ref{lattice}. As revealed in experiments, we note that the 2D layers are not stacked exactly one above the other but a finite shear displacement is observed alternatively among the adjacent layers (see Fig.~\ref{lattice}(a)-(c)). However, as the effective region  of interlayer interaction is associated with the H-bonds, we, therefore, consider the $J_3$ interaction to be of short-range (adjacent layers) composed of simple two-spin Ising interactions. We here consider the coupling $J_0$ as the largest among all the couplings  until stated explicitly, for the subsequent analysis carried throughout this work. We fix the value of $J_0$ equal to 1.  The four competing local interactions in Eq.~\eqref{Hamiltonian} except the Zeeman term (represented by $K$) are associated with different proton dynamics~\cite{vijigiri2018classical}. \\
\indent
Now we briefly describe the summary of previous work carried in the absence of the $J_3$ term. We request the interested readers to go through the Ref.~\onlinecite{vijigiri2020dipole,vijigiri2018classical} for more details. The four-body interaction term given by $J_0$ has eight degenerate spin configurations for a given plaquette. Among these eight, there are two configurations with all spins in 1/2 or all spins in -1/2 state and six configurations satisfying  \textit{ice-rule} where two spins are in 1/2 and rest of the other two spins in -1/2 state.  And in these six ice-rule states, there are four states that satisfy ice-rule and also have finite molecular polarization and the rest of the two states that do not have any finite polarization. It may be noted that the model described by the strong coupling term, $J_0$ and the external field $K_x$ is nothing but the quantum Ising gauge theory in (2+1) dimensions~\cite{chern2014gauge}. This theory is well-known to have a local $\mathcal{Z}_2$ symmetry where the degeneracy scales exactly equal to $2^{N_x+N_y+1}$~\cite{vijigiri2018classical}. The model shows a deconfinement to confinement transition as $K$ is increased from zero.  It is know that the addition of Ising type intramolecular interaction given by $J_1$ extends the deconfined phase further \cite{vijigiri2018classical}. The $J_2$ term in Eq.~\eqref{Hamiltonian} denotes a dipole-dipole interactions between two sub-plaquettes $A$ and $B$ as shown in Fig.~\ref{lattice}. This dipole-dipole interaction yields a global stripe ordering with four-fold degeneracy~\cite{vijigiri2018classical}. \\
\indent In an earlier study~\cite{vijigiri2020dipole}, it was found that in the presence of  dipole-dipole interaction, the Hamiltonian in Eq. \ref{Hamiltonian} with $J_3=0$ yields three different phases as the temperature (or field) is increased.  At low temperatures,   there exists an antiferroelectric phase up to a critical temperature. We note that in this global ordered antiferroelectric phase the ice-rules constraint is still satisfied but has finite dipole moment in order to minimize the $J_2$ interaction. As the temperature $T$ is further increased we find a quantum paraelectric phase (or the quantum liquid-like phase) where the long range order of antiferroelectric phase does not exist but the ice-rule configuration still survives locally maintaining a quantum coherence at short length scale. Finally as the temperature $T$, is increased further the thermal fluctuations prevail over the system and eventually the paraelectric phase is seen.  \\
\indent
In the absence of $J_3$, the spin configurations  within the two-dimensional layers are independent of each other (layers). However, when $J_3$ is turned on, a certain correlation between the spin configurations of  nearest neighbour or adjacent layers may appear.  If one considers two isolated plaquettes which are connected vertically by the $J_3$ coupling, then one can naively understand that for either FM or AFM coupling the configurations are just aligned or anti-aligned to each other respectively. On the other hand if $J_3$ is disordered there may be more choices indicating increased entropy. However this isolated picture of  two plaquettes coupled vertically can not be generalized for two dimensional layers connected by $J_3$ bonds. %It is true that a given spin-configurations at a given plaquette can be repeated or anti-aligned depending on the nature of $J_3$, the dynamics gets consider quantum as well as thermal fluctuations.
For example there are many degenerate configurations satisfying ice-rule configurations for a given two-dimensional layer. And when the layer goes from one degenerate configuration to another, this introduces certain length and time scale for the spins in adjacent layers to adjust. Thus this sets in some kind of rigidity against quantum and thermal fluctuations which could be the reason of high $T_c$ in experiments. In this respect present study seeks to quantitatively determine the dynamics of such coupled systems and how the quantum fluctuation and thermal fluctuation guides the systems to a new description of three dimensional squaric acid system.
%Intuitively, we understand that in general a) if the $J_3$ is FM, one would hope for a replica of configurations through all the layers. But for b) if the $J_3$ is AFM in nature, the layers could be trivially antiferroelectrically stacked. However, the effect on the intermediate state is not subtle, as it is purely a quantum state. And, for c) $J_3$ being disordered, one expects destruction of phases regardless of whether the state being globally ordered or a quantum superposition one. %\mandal{More on the physics due to this mechanism and a cartoon figure and then below paragraph will be completed. The cartoon figure will try to show constraints on ice-rules/dipole-dipole interaction due to AFM/FM/disordered strength of $J_3$.}{\color{blue}{This part is done: VIKAS,  updated figure, and text.}} \\
%%%%%%%%%%%%%%%%%%%%%%%%%%%%%%%%%%%%%%%%%%%%%%%
%\indent On the whole, to briefly outline the role of $J_3$, we note that the effect of it on the IR states within the 2D layers strictly depends on both the nature of the intercoupling ($J_3$) considered and also on the strength of it compared to the strengths of intralayer interactions ($J_{0,1,2}$). For example, in the AFM case, due to its frustrated tussle with $J_0$, we, therefore, believe that the entropy reduction near the intermediate liquid-like phase could be manifest (see Fig.~\ref{order_param}). Nevertheless, we proceed on to details of the simulation and the method used as briefed in the next section~\ref{method}.  \mandal{More on the role of $J_3$..what non-triviality it can bring. One thing is clear that when $J_3$ is much larger than anyone else that is for $J_3 \infty$ limit, the ice-rule should break down. Is that transition is of the first order or second order, is it possible to realize it in QMC. if yes very good otherwise, can we classically give an Adhoc estimate of that? One thing is sure, the possible ice-rule configuration in one layer now be coupled with the other layer..so the number of configurations will change...we need to estimate this. It means that entropy will be reduced. When there no coupling the entropy is higher than when the coupling is there. Can we define some parameters that can define this and in the abstract also we can write the results on this entropy change? External pressure can increase the value of $J_3$. can Paulson be of any help in doing a little bit of spin-wave calculation? It will be easy and doable within a month some primary results will be there.}  \color{blue}{Edited accordingly: VIKAS} \color{black}
%%%%%%%%%%%%%%%%%%%%%%%%%%%%%%%%%%%%%%%%%%%%%%%%%%%%%%%%%%%%%%%%%
                % Method
%%%%%%%%%%%%%%%%%%%%%%%%%%%%%%%%%%%%%%%%%%%%%%%%%%%%%%%%%%%%%%%%%
\subsection{Quantum Monte Carlo}
\label{method}
Here we briefly outline the method and detail implementation  of the Imaginary-time quantum Monte Carlo technique. We use the discretized version of the continuous time QMC algorithm developed by $Nakamura$~\cite{nakamura2008efficient}. Using the Suzuki-Trotter decomposition scheme~\cite{SUZUKI1993232}, the present three-dimensional quantum Hamiltonian given in Eq.~\eqref{Hamiltonian} is mapped on to an effective four-dimensional classical Hamiltonian where the additional fourth dimension is extended along the Imaginary-time (Trotter) axis $\tau$.  The effective classical action thus obtained is given by,
\begin{eqnarray}
\label{action}
\mathcal{S} &=& \sum_{p}\bigg(-\frac{\beta J_0}{M}\sum_{\square} A_{\square,p} + \frac{\beta J_1}{M}\sum_{\square}B_{\square,p} \nonumber \\
&& - \frac{\beta J_2}{M}\sum_{\langle AB \rangle}\vec{P}_{A_p}\cdot\vec{P}_{B_p} + \frac{\beta J_3}{M}\sum_{\langle L_p\rangle}\sigma_{L,p}\sigma_{L+1,p} \nonumber  \\
&& - K^{\prime}\sum_{i}\sigma_{i,p}\sigma_{i,p+1} \bigg),
\end{eqnarray}
%Studying the finite temperature property of an interacting spin system especially with a four-spin interacting Hamiltonian is not a  straightforward task. The analytical methods which rely on fermionization of the spins which map the original spin problem into interacting fermions face limitations of approximation methods and for such a spin system which contains a four-spin interaction usually results in six or eight fermion interacting models, thus where the analytical method became more intractable.
where at each Trotter index $p$ the operators initially of the form $\hat{A}_{\square}$ are mapped to $A_{\square,p}$. Note that the quantum spins are now replaced by their classical Ising variables, for instance,  $A_{\square,p}=\sigma_{1,p}\sigma_{2,p}\sigma_{3,p}\sigma_{4,p}$ and similarly the same is applied for all the remaining operators in the original Hamiltonian (Eq.~\eqref{Hamiltonian}). Here, $M$ is the cut-off length along Trotter axes and is fixed at the start of the simulation, while $\beta$ is the inverse temperature and $K^{\prime}=-\frac{1}{2}\ln \coth(\beta K/M)$ is the effective coupling (FM) along the Trotter dimension $\tau$.\\% The operators are now replaced by their classical variables of spin $\sigma_{i,p}$ with $i$ being the site-index and $p$ the imaginary-time.  \\
\indent We adapt the Swendsen-Wang (SW) \cite{swendsen1987nonuniversal} cluster update applied along the imaginary-time direction and interlayer axes. However, the case for $J_2=0$ is highly degenerate at smaller fields, and to avoid apparent spin-freezing and to speed up the relaxation process we employ the variant of the loop-flip update algorithm~\cite{Rahman}. Our variant of the algorithm involves loop that are formed at a particular imaginary-time slice randomly and extending them as SW clusters along the Trotter and interlayer axes. Once all such clusters are identified we update the configurations with their respective heat-bath probabilities. We note that for the case of $J_2\neq 0$ the loops are flipped with probability proportional to their Boltzmann weight since they are no more iso-energetic even though they form the subset of the restricted ice-rules. We consider system sizes up to $N=12\times12\times8$ block of spatial spins comprising the 3D lattice and the number of Trotter slices are fixed up to $M=1000$. About $\sim 10^5$ Monte Carlo steps(MCS) were used for equilibration and $5\times10^5$ were used for measurements. Binning analyses were done by dividing results into six bins where statistical errors are estimated. \\
\indent As we are interested in analyzing the system under both the quantum (fields) and thermal (temperature) fluctuations we, therefore, vary both the parameters taken along the contour that is parametrized by $\theta$, such that $\theta=\tan^{-1}(K\beta)$. For small inverse temperatures, the system is subjected to strong thermal fluctuations and vice-versa. We now proceed on to the next section~\ref{method} describing the details of physical quantities that will be estimated using the QMC method briefed here.
%%%%%%%%%%%%%%%%%%%%%%%%%%%%%%%%%%%%%%%%%%%%%%%%%%%%%%%%%%%%%%%%%
                % Physical Quantities
%%%%%%%%%%%%%%%%%%%%%%%%%%%%%%%%%%%%%%%%%%%%%%%%%%%%%%%%%%%%%%%%%
